\documentclass{beamer}
\usepackage[utf8]{inputenc}

\usetheme{Madrid}
\usecolortheme{default}
\usepackage{amsmath,amssymb,amsfonts,amsthm}
\usepackage{txfonts}
\usepackage{tkz-euclide}
\usepackage{listings}
\usepackage{adjustbox}
\usepackage{array}
\usepackage{tabularx}
\usepackage{gvv}
\usepackage{lmodern}
\usepackage{circuitikz}
\usepackage{tikz}
\usepackage[utf8]{inputenc}

\usetheme{Madrid}
\usecolortheme{default}
\usepackage{amsmath,amssymb,amsfonts,amsthm}
\usepackage{txfonts}
\usepackage{tkz-euclide}
\usepackage{listings}
\usepackage{adjustbox}
\usepackage{array}
\usepackage{tabularx}
\usepackage{gvv}
\usepackage{lmodern}
\usepackage{circuitikz}
\usepackage{tikz}
\usepackage{graphicx}

\setbeamertemplate{page number in head/foot}[totalframenumber]

\usepackage{tcolorbox}
\tcbuselibrary{minted,breakable,xparse,skins}



\definecolor{bg}{gray}{0.95}
\DeclareTCBListing{mintedbox}{O{}m!O{}}{%
  breakable=true,
  listing engine=minted,
  listing only,
  minted language=#2,
  minted style=default,
  minted options={%
    linenos,
    gobble=0,
    breaklines=true,
    breakafter=,,
    fontsize=\small,
    numbersep=8pt,
    #1},
  boxsep=0pt,
  left skip=0pt,
  right skip=0pt,
  left=25pt,
  right=0pt,
  top=3pt,
  bottom=3pt,
  arc=5pt,
  leftrule=0pt,
  rightrule=0pt,
  bottomrule=2pt,

  colback=bg,
  colframe=orange!70,
  enhanced,
  overlay={%
    \begin{tcbclipinterior}
    \fill[orange!20!white] (frame.south west) rectangle ([xshift=20pt]frame.north west);
    \end{tcbclipinterior}},
  #3,
}
\lstset{
    language=C,
    basicstyle=\ttfamily\small,
    keywordstyle=\color{blue},
    stringstyle=\color{orange},
    commentstyle=\color{green!60!black},
    numbers=left,
    numberstyle=\tiny\color{gray},
    breaklines=true,
    showstringspaces=false,
}
%------------------------------------------------------------
%This block of code defines the information to appear in the
%Title page
\title %optional
{5.4.9}
\date{October  2025}
%\subtitle{A short story}

\author % (optional)
{Bhoomika L - EE25BTECH11014}

\begin{document}
\frame{\titlepage}
\begin{frame}{Question}Using elementary transformations, find the inverse of the following matrix\\
\centering$\myvec{1&2\\3&4}$
\end{frame}

\begin{frame}{Solution}
We know that,
\begin{align}
     \vec{A^{-1}}\vec{A}=\vec{I}
\end{align}
     where $\vec{I}$ is a 2 X 2 identity matrix.\\
      The following is the augmented matrix we get,
      \begin{align}
           \myvec{
    \begin{array}{cc|cc}
      1 & 2 & 1 & 0 \\ 
      3 & 4 & 0 & 1
    \end{array}
  } \end{align}
\end{frame}

\begin{frame}{Solution}
\begin{align}
     \myvec{
    \begin{array}{cc|cc}
      1 & 2 & 1 & 0 \\ 
      3 & 4 & 0 & 1
    \end{array}                    
  }                                 
\xrightarrow[R_2 \longrightarrow \frac{R_2}{-2}]{R_2 \longrightarrow R_2 - 3 R_1}  \myvec{
    \begin{array}{cc|cc}
      1 & 2 & 1 & 0 \\ 
      0 & 1 & \frac{3}{2} & \frac{-1}{2}
    \end{array}
  }\\
   \myvec{\begin{array}{cc|cc}
      1 & 2 & 1 & 0 \\ 
      0 & 1 & \frac{3}{2} & \frac{-1}{2}
    \end{array}
  }\xrightarrow{\;R_1\longrightarrow R_1-2R_2\;}   \myvec{\begin{array}{cc|cc}
      1 & 0 & -2 & 1 \\ 
      0 & 1 & \frac{3}{2} & \frac{-1}{2}
 \end{array}}\\
 Therefore,\vec{A^{-1}}=\frac{1}{2}\myvec{-4 & 2 \\3 & -1}.
\end{align}
\end{frame}

 \begin{frame}[fragile]
\frametitle{C Code }
\begin{lstlisting}
#include <stdio.h>

int inverse(double* mat, double* result) {
    double a = mat[0], b = mat[1], c = mat[2], d = mat[3];
    double det = a*d - b*c;
    if(det == 0) return 0;
    
    result[0] =  d/det;
    result[1] = -b/det;
    result[2] = -c/det;
    result[3] =  a/det;
    return 1;
}

\end{lstlisting}
\end{frame}

 \begin{frame}[fragile]
\frametitle{Python Code }
\begin{lstlisting}
import ctypes

lib = ctypes.CDLL("./libinverse.so")
lib.inverse.argtypes = [ctypes.POINTER(ctypes.c_double),
                        ctypes.POINTER(ctypes.c_double)]
lib.inverse.restype = ctypes.c_int

# Input each element individually
a11 = float(input("Enter mat[1][1]: "))
a12 = float(input("Enter mat[1][2]: "))
a21 = float(input("Enter mat[2][1]: "))
a22 = float(input("Enter mat[2][2]: "))
\end{lstlisting}
\end{frame}

 \begin{frame}[fragile]
\frametitle{Python Code }
\begin{lstlisting}
# Create ctypes array row-wise
mat_array = (ctypes.c_double * 4)(a11, a12, a21, a22)
res_array = (ctypes.c_double * 4)()

# Call C function
status = lib.inverse(mat_array, res_array)

if status == 0:
    print("Matrix is singular, inverse does not exist.")
else:
    inv_matrix = [[res_array[0], res_array[1]],
                  [res_array[2], res_array[3]]]
    print("Inverse matrix:")
    for row in inv_matrix:
        print(row)
\end{lstlisting}
\end{frame}


\end{document}
